\cotefile{157-5}
\scanpage{357}\mspage{\deleted{\unsure{3|2}}\unsure{2|3}73}
% slanted
\margin{sous réserve de validité de 4.11.5.}
\uline{Corollaire} 4.11.7. Pour que $M=\unsure{\mathrm{Ind}|\mathrm I\!-\!d}_{\unsure{\mathcal F|\tau}} C$
soit stable par $\underrightarrow{\lim}$ \unsure{lorsque|(resp.} par $\underleftarrow{\lim}$ $J$ \uline{fini},
il \uline{suffit} que $C$ le soit. Et c'est
aussi nécessaire dans le cas envisagé,
lorsqu'on suppose $C$ \unsure{karoubienne}.

4.12. \uline{\deleted{\unsure{Catégories|Remontées|Foncteurs}}\inserted{Foncteurs} accessibles entre catégories
accessibles}. \deleted{\unsure{Sous-catégories}, \inserted{\deleted{\unsure{Par|P}}\unsure{Fam. d'objets|Parties d'objets},} \uline{Accessibilité d'une
catégorie accessible}.}

% slanted
\margin{\unsure{Remonter ceci} à 4.8.}
\uline{Proposition} 4.12.1. Soient $M,N$ deux
catégories accessibles\inserted{ (4.8.)}, et $f:M\unsure{-|\longrightarrow}N$
un foncteur accessible. Alors
il existe \unsure{des|deux} cardinaux \deleted{$\unsure{\pi_1}$}\inserted{infinis $\pi_0,\pi_1$} tels que
\inserted{$\pi_1\geq\pi_0$ et que} pour \deleted{\unsure{tout}} cardinal $\pi\geq\pi_1$, \deleted{\illegible{1}}
on ait
\[
 f(M_\pi)\subset N_{\pi^{\pi_0}}, \tag*{\unsure{(4.12.1.1|(4.12.1.1)}}
\]
et en particulier
\[
 f(M_\pi)\subset N_\pi\qquad\text{si }\pi=\pi^{\pi_0}. \tag*{(4.12.1.2)}
\]
De façon \inserted{plus} précise, on peut prendre
pour $\pi_0$ un cardinal \deleted{\unsure{tel}} utile pour $M$
% archive sheet number: 356
